\chapter{Описание фаззинг-тестирования} \label{ch1}

Данная глава посвящена динамическому анализу приложения с использованием
фаззинг-тестирования. Описывается выбранный фаззер, команды запуска, окружение,
а также тестовые программы для двух модулей с внесёнными ошибками.

\section{Описание фаззера и основных команд} \label{ch1:sec1}

В задании для Java указан Jazzer. Так как приложение реализовано на Node.js, для
динамического анализа использован Jazzer.js --- coverage-guided in-process fuzzing
engine для JavaScript-кода [1]. Jazzer.js основан на подходе libFuzzer
и генерирует входные данные, стремясь увеличить покрытие исполняемых ветвей
программы.

В проект добавлена зависимость:
\begin{center}
\texttt{@jazzer.js/core}
\end{center}

Основные команды фаззинга приведены в \taref{tab:fuzz-commands}.

\noindent
\begingroup
\centering
\small
\begin{longtable}[c]{|p{6cm}|p{9cm}|}
    \caption{Основные команды фаззинг-тестирования}
    \label{tab:fuzz-commands}
    \\
    \hline
    Команда & Смысл \\ \hline
    \endfirsthead
    \hline
    Команда & Смысл \\ \hline
    \endhead
    \hline
    \texttt{npm install} & Установка зависимостей проекта, включая Jazzer.js \\ \hline
    \texttt{npm run fuzz:dataset} & Запуск фаззинга модуля разбора имени датасета \\ \hline
    \texttt{npm run fuzz:percent} & Запуск фаззинга модуля декодирования percent-encoded строки \\ \hline
    \texttt{jazzer <target> --sync -- -runs=20000} & Запуск синхронного fuzz target с ограничением числа запусков \\ \hline
\end{longtable}
\normalsize
\endgroup

Jazzer.js был выбран по двум причинам. Во-первых, он является специализированным
инструментом фаззинга для Node.js, поэтому позволяет анализировать проект без
переписывания кода на Java. Во-вторых, он поддерживает привычную для Jazzer модель
fuzz target: тестируемая функция получает массив байтов, преобразует его во входной
формат приложения и передаёт в целевой модуль. Для данной работы это важно, так как
оба анализируемых дефекта связаны именно с непредсказуемым пользовательским
вводом.

Смысл фаззинг-тестирования в данной работе состоит не в проверке заранее известных
примеров, а в автоматическом подборе входов, которые нарушают устойчивость
программы. Поэтому успешным результатом считается не прохождение всех запусков, а
обнаружение минимального входа, приводящего к необработанному исключению.

Параметр \texttt{--sync} используется, так как fuzz target вызывает синхронный
JavaScript-код. Параметр \texttt{-runs=20000} ограничивает число сгенерированных
входов. При обнаружении необработанного исключения Jazzer.js завершает запуск и
сохраняет минимальный вход в файл \texttt{crash-*}.

\section{Последовательность действий по запуску фаззера и настройка окружения} \label{ch1:sec2}

Запуск фаззера выполнялся в каталоге приложения:
\begin{center}
\texttt{C:\textbackslash mtpo-lab2\textbackslash practice3\textbackslash palindrome\_app}
\end{center}

Последовательность действий:
\begin{enumerate}
    \item установить зависимости командой \texttt{npm install};
    \item подготовить fuzz target в каталоге \texttt{fuzz};
    \item добавить npm-скрипты \texttt{fuzz:dataset} и \texttt{fuzz:percent};
    \item выполнить \texttt{npm run fuzz:dataset};
    \item выполнить \texttt{npm run fuzz:percent};
    \item сохранить вывод фаззера в каталог \texttt{reports/fuzz}.
\end{enumerate}

Для исключения влияния старых результатов перед повторным запуском удалялись
файлы \texttt{crash-*}. Если фаззер снова создавал такой файл, это означало, что
ошибка воспроизводится. Отчёты сохранялись в текстовом виде и затем включались в
приложение к отчёту. После сохранения вывод был перекодирован в UTF-8, так как
PowerShell может записывать перенаправленный вывод в UTF-16LE, что приводит к
ошибкам LaTeX при подключении листингов.

Программная конфигурация приведена в \taref{tab:fuzz-env}.

\noindent
\begingroup
\centering
\small
\begin{longtable}[c]{|p{5cm}|p{10cm}|}
    \caption{Окружение фаззинг-тестирования}
    \label{tab:fuzz-env}
    \\
    \hline
    Компонент & Значение \\ \hline
    \endfirsthead
    \hline
    Компонент & Значение \\ \hline
    \endhead
    \hline
    Операционная система & Windows \\ \hline
    Среда выполнения & Node.js v24.13.0 \\ \hline
    npm & 11.6.2 \\ \hline
    Фаззер & Jazzer.js, пакет \texttt{@jazzer.js/core} \\ \hline
    Каталог отчётов & \texttt{reports/fuzz} \\ \hline
\end{longtable}
\normalsize
\endgroup

\section{Описание тестовых программ} \label{ch1:sec3}

Для фаззинга выбраны два модуля, в которые внесены две ошибки.

\textbf{Модуль \texttt{datasetNameParser.js}.}
Модуль разбирает имя датасета по формату \texttt{TYPE\_SIZE\_SERIES}. Ошибка
заключается в отсутствии проверки числа частей после \texttt{split("\_")}. Если
входная строка не содержит третьей части, выражение \texttt{parts[2].toLowerCase()}
вызывает \texttt{TypeError}.

В нормальном сценарии функция должна принимать имена вида
\texttt{UNI\_16000\_A}, \texttt{ALT\_5000\_A}, \texttt{MIR\_10000\_A}. Внесённая
ошибка проявляется на неполных именах: пустая строка, \texttt{UNI}, \texttt{UNI\_16000}.
Такие входы возможны при ручном вводе имени датасета или при повреждении файла со
списком наборов.

Fuzz target:
\begin{lstlisting}[language=JavaScript,breaklines=true]
const { parseDatasetName } = require("../src/vulnerable/datasetNameParser");

module.exports.fuzz = function fuzz(data) {
  parseDatasetName(data.toString("utf8"));
};
\end{lstlisting}

Ожидаемое безопасное поведение для этого модуля --- выброс контролируемого
исключения с сообщением о неверном формате имени датасета. Фактическое поведение
--- необработанный \texttt{TypeError}, указывающий на ошибку валидации.

\textbf{Модуль \texttt{percentInputDecoder.js}.}
Модуль декодирует percent-encoded строку через \texttt{decodeURIComponent}. Ошибка
заключается в отсутствии обработки исключения \texttt{URIError}, возникающего на
некорректной последовательности percent-кодирования, например на строке
\texttt{\%}.

В нормальном сценарии функция может получать строки вида \texttt{abacaba},
\texttt{a\%20b\%20a} или \texttt{\%D0\%B0}. Некорректный вход \texttt{\%} является
частичным percent-encoded фрагментом. Если такая строка поступит от пользователя,
приложение должно выдать диагностическое сообщение, а не завершать обработку
аварийным исключением.

Fuzz target:
\begin{lstlisting}[language=JavaScript,breaklines=true]
const { decodePercentInput } = require("../src/vulnerable/percentInputDecoder");

module.exports.fuzz = function fuzz(data) {
  decodePercentInput(data.toString("utf8"));
};
\end{lstlisting}

Ожидаемое безопасное поведение для этого модуля --- обработать \texttt{URIError}
и вернуть понятную ошибку входного формата. Фактическое поведение --- исключение
пробрасывается наружу.

\section{Классификация тестовых программ} \label{ch1:sec4}

Сводная информация о выбранных тестовых программах приведена в
\taref{tab:fuzz-programs}.

\noindent
\begingroup
\centering
\small
\begin{longtable}[c]{|p{4cm}|p{3.5cm}|p{3.5cm}|p{4cm}|}
    \caption{Тестовые программы для фаззинга}
    \label{tab:fuzz-programs}
    \\
    \hline
    Fuzz target & Тестируемый модуль & Тип входа & Ожидаемый класс дефекта \\ \hline
    \endfirsthead
    \hline
    Fuzz target & Тестируемый модуль & Тип входа & Ожидаемый класс дефекта \\ \hline
    \endhead
    \hline
    \texttt{datasetNameFuzzTarget.js} & \texttt{datasetNameParser.js} & Произвольная строка имени датасета & Ошибка валидации структуры строки \\ \hline
    \texttt{percentDecoderFuzzTarget.js} & \texttt{percentInputDecoder.js} & Произвольная percent-encoded строка & Ошибка обработки некорректного кодирования \\ \hline
\end{longtable}
\normalsize
\endgroup

Оба fuzz target построены минимально: они не фильтруют входы и не подавляют
исключения. Такой подход выбран намеренно, потому что цель эксперимента ---
проверить устойчивость анализируемых функций к произвольному пользовательскому
вводу.

\section{Выводы} \label{ch1:conclusion}

Для динамического анализа подготовлены два fuzz target, соответствующие двум
модулям с внесёнными ошибками. Фаззер Jazzer.js выбран как аналог Jazzer для
JavaScript/Node.js и позволяет находить необработанные исключения на автоматически
сгенерированных входных данных.
